% ============================================================
% CAPITOLUL 3 — TEHNOLOGII ȘI INSTRUMENTE UTILIZATE
% ============================================================

\chapter{Tehnologii și Instrumente Utilizate}
\label{ch:tehnologii}

\section{Limbaje de programare pentru calcul științific}
\label{sec:limbaje}

Alegerea limbajului de programare pentru implementarea modelelor matematice
nu este neutră: performanța, expresivitatea și ecosistemul de biblioteci
au un impact direct asupra eficienței dezvoltării și calității rezultatelor.
Această secțiune prezintă fundamentele și caracteristicile tehnice ale
celor patru limbaje utilizate.

% ---- R ---------------------------------------------------
\subsection{R}
\label{subsec:r}

\subsubsection{Istoric și filosofie}

R este un limbaj de programare și mediu software pentru calcul statistic și
grafică, descendent direct al limbajului S (Bell Laboratories, 1976). A fost
creat de Ross Ihaka și Robert Gentleman la Universitatea din Auckland în 1993,
iar în 1997 a fost adoptat ca proiect open-source de \emph{R Development Core
Team} \cite{rcoreteam2024}. Filosofia R este centrată pe date: vectorii,
matricile și cadrele de date (data frames) sunt cetățeni de primă clasă,
iar gramatica limbajului este proiectată pentru manipularea și analiza datelor.

\subsubsection{Caracteristici tehnice relevante}

R este un limbaj \emph{interpretat}, dinamic tipizat, cu suport pentru
programare funcțională, orientată pe obiecte (sistemele S3, S4 și R5/R6)
și procedurală. Caracteristici relevante pentru calculul științific:

\begin{itemize}
    \item \textbf{Vectorizare nativă:} operațiile aritmetice și logice aplicate
    pe vectori și matrice sunt automat vectorizate și executate prin cod C/Fortran
    compilat, evitând buclele interpretate costisitoare.
    \item \textbf{Lazy evaluation:} argumentele funcțiilor sunt evaluate doar
    când sunt necesare, permițând construcții expresive.
    \item \textbf{Gestionarea memoriei:} R utilizează un garbage collector cu
    copiere la scriere (copy-on-write), care poate genera overhead semnificativ
    pentru structuri de date mari.
    \item \textbf{Interfețe C/Fortran:} prin funcțiile \texttt{.C()},
    \texttt{.Call()} și \texttt{.Fortran()}, R poate apela cod compilat nativ,
    stând la baza multor pachete de înaltă performanță.
\end{itemize}

\subsubsection{Pachete relevante pentru modelul SIR}

\begin{itemize}
    \item \textbf{deSolve} \cite{hairer1993}: pachetul principal pentru
    rezolvarea sistemelor de EDO, ecuații diferențiale cu derivate parțiale
    (EDP) și ecuații algebrico-diferențiale (EAD). Integrează solver-e
    consacrate: \texttt{lsoda} (Adams/BDF adaptiv), \texttt{ode45}
    (Dormand-Prince), \texttt{rk4} (Runge-Kutta fix). Interfața principală
    este funcția \texttt{ode()}.
    \item \textbf{ggplot2}: sistem de grafică bazat pe \emph{grammar of graphics},
    utilizat pentru vizualizarea traiectoriilor SIR.
    \item \textbf{data.table}: manipulare eficientă (în memorie) a seturilor
    de date mari, cu sintaxă concisă.
    \item \textbf{microbenchmark} și \textbf{bench}: măsurarea precisă a
    timpilor de execuție pentru benchmarking.
\end{itemize}

% ---- JULIA -----------------------------------------------
\subsection{Julia}
\label{subsec:julia}

\subsubsection{Istoric și filosofie}

Julia a fost creat la MIT de Jeff Bezanson, Alan Edelman, Stefan Karpinski
și Viral B. Shah, cu versiunea 1.0 lansată în august 2018 \cite{bezanson2017julia}.
Motivația proiectului a fost depășirea \emph{problemei celor două limbaje}
(two-language problem): cercetătorii scriu prototipuri în R/Python (productivitate
ridicată, viteză scăzută) și re-implementează codul critic în C/Fortran (viteză
ridicată, productivitate scăzută). Julia vizează ambele obiective simultan.

\subsubsection{Compilare Just-In-Time și sistemul de tipuri}

Julia utilizează compilare \textbf{JIT} (Just-In-Time) bazată pe LLVM, cu
o particularitate fundamentală: specializarea pe tipuri concrete. Când o funcție
este apelată pentru prima dată cu argumente de tipuri specifice, compilatorul
generează cod mașină optimizat pentru acele tipuri exacte. Apelurile ulterioare
cu aceleași tipuri rulează direct codul compilat, apropiind performanța de C.

\begin{itemize}
    \item \textbf{Multiple dispatch:} funcțiile pot fi definite cu comportamente
    diferite pentru combinații diferite de tipuri ale argumentelor (metode).
    Compilatorul selectează și compilează metoda corectă la momentul apelului.
    \item \textbf{Sistem de tipuri parametric:} tipurile pot fi parametrizate
    (ex.\ \texttt{Array\{Float64, 1\}}), permițând generarea eficientă de cod
    specializat.
    \item \textbf{Metaprogramming:} macrocomenzile Julia ($\texttt{@macro}$)
    permit generarea și transformarea codului la momentul parsării.
\end{itemize}

\subsubsection{DifferentialEquations.jl}

Ecosistemul \textbf{SciML} (Scientific Machine Learning) al Julia conține
pachetul \texttt{DifferentialEquations.jl} \cite{rackauckas2017diffeq}, cel
mai cuprinzător framework open-source pentru rezolvarea ecuațiilor diferențiale.
Acesta include:
\begin{itemize}
    \item Peste 300 de algoritmi de rezolvare pentru EDO, EDP, EAD,
    ecuații diferențiale stochastice (EDS) și ecuații cu întârziere (DDEs);
    \item Solver-e de înaltă performanță: \texttt{Tsit5} (Tsitouras RK 4/5),
    \texttt{Vern7/9} (Verner 7/9), \texttt{Rodas4/5} (pentru probleme stiff);
    \item Integrare nativă cu GPU-uri și calcul distribuit;
    \item Diferențiere automată (AD) pentru estimarea parametrilor și
    identificarea sistemelor.
\end{itemize}

% ---- PYTHON ----------------------------------------------
\subsection{Python}
\label{subsec:python}

\subsubsection{Istoric și filosofie}

Python a fost creat de Guido van Rossum și lansat în 1991 \cite{vanrossum2009python}.
Filosofia sa, rezumată în \emph{Zen of Python} (\texttt{PEP 20}), accentuează
lizibilitatea și simplitatea codului. Python a devenit limbajul dominant în
știința datelor și machine learning datorită ecosistemului său extins și
comunității masive.

\subsubsection{Caracteristici tehnice și ecosistem științific}

Python este un limbaj interpretat, dinamic tipizat, multi-paradigmă. Principala
limitare pentru calcul intensiv este \textbf{GIL} (Global Interpreter Lock),
care previne execuția paralelă a thread-urilor Python pure. Performanța este
recuperată prin biblioteci care rulează cod nativ C/Fortran în afara GIL:

\begin{itemize}
    \item \textbf{NumPy}: calculul numeric vectorizat pe array-uri dense;
    operațiile sunt implementate în C/Fortran și nu sunt supuse GIL.
    \item \textbf{SciPy} \cite{virtanen2020scipy}: algoritmică științifică,
    incluzând modulul \texttt{scipy.integrate} cu funcțiile \texttt{odeint}
    (wrapper LSODA) și \texttt{solve\_ivp} (interfață modernă cu multiple
    metode: RK45, DOP853, Radau, BDF, LSODA).
    \item \textbf{Matplotlib}: vizualizare 2D/3D flexibilă.
    \item \textbf{Pandas}: manipularea datelor tabulare.
\end{itemize}

\subsubsection{Funcțiile de rezolvare a EDO în SciPy}

Interfața \texttt{solve\_ivp} din \texttt{scipy.integrate} este apelată astfel:
\begin{lstlisting}[language=Python, caption={Interfața solve\_ivp din SciPy pentru rezolvarea sistemului SIR}]
from scipy.integrate import solve_ivp
import numpy as np

def sir_model(t, y, beta, gamma, N):
    S, I, R = y
    dSdt = -beta * S * I / N
    dIdt =  beta * S * I / N - gamma * I
    dRdt =  gamma * I
    return [dSdt, dIdt, dRdt]

# Parametri
N = 1_000_000; beta = 0.3; gamma = 0.1
S0, I0, R0 = N - 1, 1, 0
t_span = (0, 160); t_eval = np.linspace(0, 160, 1601)

sol = solve_ivp(
    fun=lambda t, y: sir_model(t, y, beta, gamma, N),
    t_span=t_span,
    y0=[S0, I0, R0],
    method='RK45',
    t_eval=t_eval,
    rtol=1e-8,
    atol=1e-10
)
\end{lstlisting}

% ---- SAS -------------------------------------------------
\subsection{SAS}
\label{subsec:sas}

\subsubsection{Istoric și model de licențiere}

SAS (Statistical Analysis System) a fost dezvoltat la North Carolina State
University începând cu 1966 și comercializat de SAS Institute din 1976.
Spre deosebire de R, Julia și Python, SAS este un software proprietar cu
licențiere comercială, larg utilizat în industria farmaceutică, financiară
și în cercetarea clinică, în parte datorită cerințelor reglementare (FDA
validation, CFR 21 Part 11).

\subsubsection{PROC IML pentru calcul numeric}

\textbf{PROC IML} (Interactive Matrix Language) este procedura SAS care
permite programare matriceală generală, analogă cu limbajele de tip MATLAB.
Prin IML, sistemul SIR poate fi rezolvat implementând manual schema RK4:

\begin{lstlisting}[language=SAS, caption={Rezolvarea sistemului SIR în SAS/IML cu schema RK4}]
proc iml;
  start sir_rhs(t, y, parms);
    beta = parms[1]; gamma = parms[2]; N = parms[3];
    S = y[1]; I = y[2]; R = y[3];
    dSdt = -beta * S * I / N;
    dIdt =  beta * S * I / N - gamma * I;
    dRdt =  gamma * I;
    return (dSdt || dIdt || dRdt);
  finish;

  start rk4_step(t, y, h, parms);
    k1 = h * sir_rhs(t,       y,           parms);
    k2 = h * sir_rhs(t + h/2, y + k1/2,   parms);
    k3 = h * sir_rhs(t + h/2, y + k2/2,   parms);
    k4 = h * sir_rhs(t + h,   y + k3,     parms);
    return (y + (k1 + 2*k2 + 2*k3 + k4) / 6);
  finish;

  /* Parametri si conditii initiale */
  N = 1000000; beta = 0.3; gamma = 0.1;
  parms = beta || gamma || N;
  y = (N-1) || 1 || 0;
  t = 0; h = 0.1; tf = 160;

  results = j(0, 4);
  do while (t <= tf);
    results = results // (t || y);
    y = rk4_step(t, y, h, parms);
    t = t + h;
  end;
quit;
\end{lstlisting}

\subsubsection{Comparație cu alternativele open-source}

SAS prezintă avantaje în medii de producție reglementate (trasabilitate,
audit trail, validare formală), dar dezavantaje semnificative pentru cercetare
exploratorie: cost de licențiere ridicat, ciclu de actualizare lent și
un ecosistem de pachete mult mai restrâns comparativ cu R și Python.
Varianta \textbf{SAS OnDemand for Academics} oferă acces gratuit pentru
uz educațional.

% ---- COMPARAȚIE -----------------------------------------
\subsection{Analiză comparativă a limbajelor}
\label{subsec:comparatie_limbaje}

\Cref{tab:comparatie_limbaje} sintetizează principalele caracteristici ale
celor patru limbaje din perspectiva utilizării pentru implementarea modelului SIR.

\begin{table}[htbp]
\centering
\caption{Comparație a limbajelor de programare utilizate}
\label{tab:comparatie_limbaje}
\renewcommand{\arraystretch}{1.3}
\begin{tabular}{lcccc}
\toprule
\textbf{Criteriu} & \textbf{R} & \textbf{Julia} & \textbf{Python} & \textbf{SAS} \\
\midrule
Licență          & Open-source & Open-source & Open-source & Proprietar \\
Paradigmă        & Funcț./OOP  & Multi       & Multi       & Procedural \\
Compilare        & Interpretat & JIT (LLVM)  & Interpretat & Interpretat \\
Solver EDO       & deSolve     & DiffEq.jl   & SciPy       & PROC IML  \\
Performanță rel. & Medie       & Înaltă      & Medie       & Medie     \\
Maturitate ecos. & Înaltă      & Medie       & Înaltă      & Înaltă    \\
Utilizare acad.  & Ridicată    & Medie       & Ridicată    & Ridicată  \\
\bottomrule
\end{tabular}
\end{table}

% ============================================================
\section{Baze de date distribuite — YugabyteDB}
\label{sec:yugabytedb}
% ============================================================

\subsection{Concepte de baze de date distribuite}
\label{subsec:baze_distribuite}

O \emph{bază de date distribuită} este o colecție de date distribuite fizic
pe mai multe noduri (mașini sau locații geografice), gestionată de un sistem
software care oferă transparența localizării: aplicațiile interacționează cu
baza de date ca și cum ar fi centralizată.

\begin{definitie}[Teorema CAP]
\label{def:cap}
Teorema CAP (Brewer, 2000) afirmă că un sistem distribuit nu poate garanta
simultan mai mult de două dintre următoarele trei proprietăți:
\begin{itemize}
    \item \textbf{Consistency (C):} toate nodurile văd aceleași date în același
    moment (citirile reflectă întotdeauna ultima scriere confirmată);
    \item \textbf{Availability (A):} orice cerere validă primește un răspuns
    (eventual fără ultima versiune a datelor);
    \item \textbf{Partition tolerance (P):} sistemul continuă să funcționeze
    chiar dacă comunicarea între noduri este parțial întreruptă.
\end{itemize}
\end{definitie}

Bazele de date distribuite moderne (YugabyteDB, CockroachDB, Spanner) optează
pentru \textbf{CP} (consistency + partition tolerance), sacrificând disponibilitatea
perfectă în favoarea coerenței datelor.

\subsection{Arhitectura YugabyteDB}
\label{subsec:yugabyte_arhitectura}

YugabyteDB \cite{yugabyte2024} este o bază de date distribuită open-source,
construită pentru a combina caracteristicile PostgreSQL cu scalabilitatea
orizontală și toleranța la erori a sistemelor NoSQL. Arhitectura sa se bazează
pe doi piloni principali:

\textbf{1. Stratul de stocare DocDB.} DocDB este motorul de stocare
distribuit al YugabyteDB, bazat pe arhitectura \textbf{LSM-tree} (Log-Structured
Merge-tree) derivată din \textbf{RocksDB}. Datele sunt stocate în \emph{tablete}
--- unități de sharding distribuite pe noduri. Fiecare tabletă este replicată
pe cel puțin trei noduri prin \textbf{protocolul Raft} (consens distribuit):
\begin{itemize}
    \item Un nod este ales \emph{lider}; scrierile trec prin lider;
    \item Leaderul repl\-ică modificarea pe nodurile urmașe (\emph{followers});
    \item O scriere este confirmată după ce o \emph{majoritate} (quorum) de
    noduri o confirmă.
\end{itemize}

\textbf{2. Nivelurile de interfață API.}
\begin{itemize}
    \item \textbf{YSQL} (YugabyteDB SQL): compatibil cu PostgreSQL 11+,
    suportând SQL complet, tranzacții ACID distribuite, join-uri, indecși,
    triggere și proceduri stocate;
    \item \textbf{YCQL} (YugabyteDB CQL): compatibil cu Apache Cassandra,
    optimizat pentru workload-uri de tip wide-column cu disponibilitate extinsă.
\end{itemize}

\subsection{Sharding și replicare}
\label{subsec:sharding}

YugabyteDB suportă două strategii de sharding:
\begin{itemize}
    \item \textbf{Hash sharding:} cheia primară este hașurată și distribuită
    uniform pe tablete; asigură echilibru de sarcină dar nu suportă eficient
    \emph{range scans};
    \item \textbf{Range sharding:} tabletele acoperă intervale contigue de
    valori ale cheii; eficient pentru interogări de tip range, dar poate
    produce \emph{hot spots}.
\end{itemize}

\begin{observatie}
Pentru datele epidemiologice cu interogări frecvente pe intervale de date
(serii temporale), range sharding pe coloana \texttt{timestamp} sau
\texttt{date} este strategia recomandată, împreună cu un index secundar
pe codul bolii și/sau regiune geografică.
\end{observatie}

% ============================================================
\section{Orchestrarea containerelor — Kubernetes}
\label{sec:kubernetes}
% ============================================================

\subsection{Containere și nevoia de orchestrare}
\label{subsec:containere}

Containerizarea (Docker) permite împachetarea unei aplicații împreună cu
toate dependențele sale într-o unitate portabilă și reproductibilă.
Pe măsură ce sistemele cresc în complexitate și numărul de containere
depășește câteva zeci, gestionarea manuală devine infezabilă: unde rulează
fiecare container? Cum se reporneşte dacă cade? Cum se scalează?

\textbf{Kubernetes} (K8s) \cite{burns2016kubernetes, kubernetes2024} este
platforma open-source de orchestrare a containerelor, inițial dezvoltată
la Google pe baza experienței cu sistemul intern Borg, și donată CNCF
(Cloud Native Computing Foundation) în 2014.

\subsection{Arhitectura Kubernetes}
\label{subsec:k8s_arhitectura}

Un cluster Kubernetes este compus din:

\textbf{Control Plane:}
\begin{itemize}
    \item \textbf{kube-apiserver:} componenta centrală, expune API-ul K8s
    (REST/gRPC); toate celelalte componente comunică exclusiv prin API server;
    \item \textbf{etcd:} baza de date distribuită key-value (Raft-based) care
    stochează starea întregului cluster;
    \item \textbf{kube-scheduler:} atribuie pod-uri nodurilor disponibile
    pe baza resurselor disponibile și a constrângerilor de afinitate;
    \item \textbf{kube-controller-manager:} rulează controllere care reconciliază
    starea dorită cu starea actuală (Deployment, ReplicaSet, StatefulSet etc.).
\end{itemize}

\textbf{Worker Nodes:}
\begin{itemize}
    \item \textbf{kubelet:} agent pe fiecare nod, gestionează ciclul de
    viață al pod-urilor (Pull image, start/stop/restart containere);
    \item \textbf{kube-proxy:} gestionează regulile de rețea (iptables/IPVS)
    pentru accesul la servicii;
    \item \textbf{Container runtime:} interfața CRI (containerd, CRI-O).
\end{itemize}

\subsection{Resurse Kubernetes relevante}
\label{subsec:k8s_resurse}

\begin{itemize}
    \item \textbf{Pod:} unitatea atomică de deployment; unul sau mai multe
    containere cu spațiu de rețea și stocare partajate;
    \item \textbf{Deployment:} gestionează un set de pod-uri identice;
    asigură restart automat și rolling updates;
    \item \textbf{StatefulSet:} similar Deployment, dar menține identitate
    stabilă pentru pod-uri (hostname, persistent volume); \emph{esențial}
    pentru YugabyteDB;
    \item \textbf{Service:} abstractizare pentru accesul la un set de pod-uri;
    asigură load balancing și descoperire de servicii;
    \item \textbf{PersistentVolume / PersistentVolumeClaim:} stocare persistentă
    independentă de ciclul de viață al pod-ului;
    \item \textbf{ConfigMap / Secret:} configurare și date sensibile externalizate.
\end{itemize}

% ============================================================
\section{Backend API — Go}
\label{sec:golang}
% ============================================================

\subsection{Caracteristici și filosofia limbajului Go}
\label{subsec:go_caracteristici}

Go (cunoscut și ca Golang) este un limbaj de programare open-source creat
la Google de Robert Griesemer, Rob Pike și Ken Thompson, cu prima versiune
stabilă în 2012 \cite{golangspec2024}. Go a fost proiectat pentru productivitate
la scară mare, cu accent pe simplitate, compilare rapidă și concurență structurată.

Caracteristici relevante pentru construcția unui API backend:
\begin{itemize}
    \item \textbf{Compilare nativă rapidă:} Go compilează direct la cod mașin,
    producând binare statice fără dependențe externe;
    \item \textbf{Goroutine și canale (CSP):} concurența este modelată prin
    goroutine (fire de execuție ușoare, ~2KB stack inițial) și canale (typed,
    pentru comunicare sincronă/asincronă); modelul Communicating Sequential
    Processes (CSP, Hoare 1978) este preferat față de memoria partajată cu mutex;
    \item \textbf{Interfețe structurale (duck typing):} un tip implementează
    implicit o interfață dacă definește metodele necesare, fără declarație
    explicită;
    \item \textbf{Garbage collection:} GC concurrent cu latență redusă
    ($<1$ms stop-the-world în versiunile recente);
    \item \textbf{Standard library bogată:} pachetul \texttt{net/http} suportă
    HTTP/1.1, HTTP/2 și TLS out-of-the-box, fără framework extern.
\end{itemize}

\subsection{Ecosistem pentru API-uri REST}
\label{subsec:go_api}

Deși \texttt{net/http} din biblioteca standard este suficient pentru API-uri
simple, framework-uri ca \textbf{Gin} sau \textbf{Echo} adaugă:
\begin{itemize}
    \item Router performant (radix tree pentru $\mathcal{O}(\log k)$ matching);
    \item Middleware pipeline (logging, CORS, rate limiting, autentificare JWT);
    \item Binding și validare automată a parametrilor HTTP (JSON, query params,
    path params);
    \item Generare automată a documentației OpenAPI/Swagger.
\end{itemize}

% ============================================================
\section{Frontend — Vue.js}
\label{sec:vuejs}
% ============================================================

\subsection{Arhitectura Vue 3 și Composition API}
\label{subsec:vue_arhitectura}

Vue.js \cite{vuejs2024} este un framework JavaScript progresiv pentru construirea
interfețelor utilizator, creat de Evan You în 2014. Vue 3 (2020) introduce
\textbf{Composition API}, o alternativă la Options API care permite organizarea
logicii prin funcționalitate (nu prin tip --- data, methods, computed) și
reutilizarea logicii prin \emph{composables}.

Arhitectura Vue 3 se bazează pe:
\begin{itemize}
    \item \textbf{Reactivity system:} bazat pe \texttt{Proxy} ES6, permite
    urmărirea dependențelor reactive cu granularitate fină;
    \item \textbf{Virtual DOM:} reconciliere eficientă a arborelui DOM cu
    algoritmul de diferențiere O(n);
    \item \textbf{Single File Components (SFC):} fișiere \texttt{.vue}
    care încapsulează template (HTML), script (JS/TS) și style (CSS)
    într-o unitate coezivă.
\end{itemize}

\subsection{Ecosistem și biblioteci utilizate}
\label{subsec:vue_ecosistem}

\begin{itemize}
    \item \textbf{Vue Router:} rutare client-side cu suport pentru rute
    dinamice, navigare programatică și guards de autentificare;
    \item \textbf{Pinia:} state management bazat pe Composition API, mai
    simplu decât Vuex și mai bine integrat cu TypeScript;
    \item \textbf{Axios:} client HTTP cu interceptori, pentru apelurile
    REST la API-ul Go;
    \item \textbf{ECharts / Chart.js:} vizualizare date epidemiologice
    (grafice SIR, comparații de performanță, hărți de densitate).
\end{itemize}
